
\begin{tikzpicture}[node distance=3.2cm and 3cm]

\node[state] (released) {RELEASED};
\node[state,right=of released] (pressed) {PRESSED};

\draw[flow] (released) to[bend left=22]
    node[above,font=\footnotesize,align=center]
    {pressione / richiesta} (pressed);

\draw[flow] (pressed) to[bend left=22]
    node[below,font=\footnotesize]
    {rilascio} (released);

\node[state,below=of released] (disabled) {DISABLED};
\node[state,right=of disabled] (active) {ACTIVE};

\draw[flow] (disabled) to[bend left=22]
    node[above,font=\footnotesize,align=center]
    {MANUAL confermato /\\ acquisizione riferimento} (active);

\draw[flow] (active) to[bend left=22]
    node[below,font=\footnotesize]
    {uscita da MANUAL} (disabled);

\end{tikzpicture}

\par\smallskip

\begin{minipage}{0.94\linewidth}
\footnotesize
\raggedright%

\textbf{Condizioni operative:}

\begin{itemize}
    \setlength{\itemsep}{2pt}
    \setlength{\parsep}{0pt}
    \setlength{\topsep}{3pt}

    \item \textbf{Pulsante:}
    una richiesta per pressione, senza ripetizioni fino al rilascio.

    \item \textbf{Potenziometro:}
    in ACTIVE invia una richiesta per variazioni di almeno 5 punti percentuali o al raggiungimento di un estremo con variazione non nulla.

    \item \textbf{Stati iniziali:}
    RELEASED e DISABLED.
\end{itemize}

\end{minipage}
